Let $R = \sqrt{\dfrac{100}{\pi}}$ be the pen radius and let the rope length be $tR$, $0<t<2$. The grazing region is the overlap of circles with radii $R,tR$ whose centres are $R$ apart. Adding the two sectors and subtracting their two triangles gives the overlap area divided by $R^2$ as $G (t) = t^2 \arccos (\dfrac{t}{2}) + \arccos (1 - \dfrac{t^2}{2}) - \dfrac{t}{2} \sqrt{4-t^2}$. The target is $G (t) = \dfrac{\pi}{2}$. The accessible region grows strictly with rope length on this interval, so there is at most one solution; continuity and the endpoint areas ensure one exists. Bisection gives $t\approx1.158728473$, so the rope length is $t \sqrt{\dfrac{100}{\pi}} \approx 6.537425346$ feet, or $\boxed{6.537\text{ feet}}$.
\begin{asy}
import graph;
usepackage("amsmath");
import geometry;
size(7cm);
pair O=(0,0),A=(1,0); real t=1.158728473; pair P=(1-t*t/2,sqrt(1-(1-t*t/2)^2));draw(circle(O,1));draw(circle(A,t),dashed);draw(O--A--P--O);label("$O$",O,(-1,-1));label("tie point $A$",A,(1,-1));label("$P$",P,(0,1));label("$R$",O/2+P/2,(-1,0));label("rope",(A+P)/2,(1,1));
\end{asy}

